% Generated from THESIS_DRAFT.md.
\chapter{Kết luận}

\section{Thành quả đã đạt được}

Khóa luận đã xây dựng một hướng bổ sung bộ nhớ sự kiện thị giác cho SmolVLA trong tác vụ nhấc và đặt lại bát. Hệ thống liên kết detection với đĩa gốc, ghi nhận departure/return bằng máy trạng thái và đưa snapshot đã kiểm nguồn vào prompt. Các cơ chế thời điểm khả dụng, incomplete history và giới hạn token giúp biểu đạt rõ thông tin nào được sử dụng tại mỗi quyết định.

Một quy trình thích ứng action expert bằng LoRA đã được triển khai và tạo ra adapter sau 3.000 bước huấn luyện. Các file export được xác minh theo manifest. Đợt rollout B0/F1 trên T1/T3 đã tạo đủ artifact cho 40 task-episodes, cho phép đối chiếu kết quả với prompt và trạng thái memory. Đây là bằng chứng về việc phương pháp đã được đưa vào một hệ thống thực thi, vượt qua mức chỉ có kiến trúc đề xuất.

\section{Điều rút ra từ kết quả hiện tại}

Kết quả hiện tại chưa chứng minh mục tiêu nâng cao hiệu năng đã đạt. F1 có một success trong 20 task-episodes trong khi B0 không có success, nhưng phép so sánh còn gộp adaptation và memory, số nhóm state nhỏ và chỉ có một seed. Không có cơ sở chuyển các con số này thành một cải thiện chính thức trên LIBERO-Mem hoặc một kết luận tổng quát về memory cho VLA.

Một kết luận có căn cứ hơn là chất lượng và khả năng sử dụng thông tin memory phải được kiểm tra độc lập với loss huấn luyện. Hầu hết snapshot cuối của F1 báo lịch sử chưa đầy đủ; success duy nhất không đi kèm việc writer đếm đủ các chu kỳ được evaluator ghi nhận. Vì vậy, việc đã tạo prompt, đã giảm loss offline và đã có rollout thành công là ba bằng chứng khác nhau, không thể thay thế cho nhau trong lập luận về hiệu quả của phương pháp.

Giá trị học thuật hiện có nằm ở cách đặt và làm rõ bài toán, thiết kế biểu diễn sự kiện có truy vết, cùng những phép kiểm giúp tránh kết luận vượt bằng chứng. Giá trị này cần được đánh giá trong phạm vi một khóa luận xây dựng và khảo sát hệ thống; chưa phải chứng minh tính mới toàn cầu hay ưu thế so với các kiến trúc memory lớn.

\section{Hạn chế và hướng phát triển}

Các hạn chế chính là phạm vi một quan hệ bát–đĩa, thiếu đánh giá event fidelity, chưa có F0 matched và retention của adapter hiện tại, cùng khác biệt chưa được đóng hoàn toàn giữa dữ liệu tác giả và runtime local. Chuỗi tạo dữ liệu derived mới còn cần export thật và audit hoàn chỉnh trước khi được dùng cho kết luận huấn luyện mới.

Hướng tiếp theo là hoàn thiện chuỗi dữ liệu có version, xây dựng đối chứng thích ứng không memory trên cùng demonstration, đo chất lượng writer và thực hiện paired evaluation với phạm vi holdout được công bố rõ. Sau khi làm rõ giá trị của lịch sử trong bài toán hẹp, nghiên cứu mới có cơ sở mở rộng sang nhiều thực thể, nhiều backbone hoặc robot thật. Tóm tắt và kết luận cuối cùng sẽ được cập nhật theo kết quả thực nghiệm bổ sung, kể cả khi kết quả đó không xác nhận giả thuyết ban đầu.
